\chapter{Methodology}

\section{One-Class Graph Neural Networks}
\label{sec:ocgin}

One-Class Classification (OCC) aims to distinguish normal data from anomalies using only normal training examples. When applied to graph data, the challenge is to learn a representation that captures the structural and attributive patterns of normal graphs or nodes. We adopt the \ac{OCGIN} architecture \citep{wang2021ocgin} as the base model, which combines \acp{GIN} \citep{xu2018gin} with the Deep \ac{SVDD} objective \citep{ruff2018deepsvdd}.

\subsection{Graph Isomorphism Network Encoder}
\label{sec:gin}

For a graph $G = (V, E)$ with node features $\mathbf{X} \in \mathbb{R}^{|V| \times d_{\text{in}}}$, a \ac{GIN} produces node embeddings through iterative neighborhood aggregation. The $k$-th layer computes

\begin{equation}
\mathbf{h}_v^{(k+1)} = \mathrm{MLP}^{(k)}\!\left( (1 + \varepsilon) \cdot \mathbf{h}_v^{(k)} + \sum_{u \in \mathcal{N}(v)} \mathbf{h}_u^{(k)} \right),
\label{eq:gin}
\end{equation}

where $\mathbf{h}_v^{(0)} = \mathbf{x}_v$ are the initial node features, $\mathcal{N}(v)$ denotes the set of neighbors of $v$, $\varepsilon$ is a learnable or fixed parameter, and $\mathrm{MLP}^{(k)}$ is a multi-layer perceptron consisting of two linear layers with a ReLU activation in between. The GIN architecture is theoretically maximally expressive among message-passing \acp{GNN} under the injective aggregation condition \citep{xu2018gin}.

After $L$ layers, we obtain final node embeddings $\mathbf{Z} \in \mathbb{R}^{|V| \times d}$ where $\mathbf{z}_v = \mathbf{h}_v^{(L)}$.

\subsection{Graph-Level Pooling}
\label{sec:pooling}

For graph-level tasks, the node embeddings are aggregated into a single graph embedding via global mean pooling:

\begin{equation}
\mathbf{z}_G = \frac{1}{|V_G|} \sum_{v \in V_G} \mathbf{h}_v^{(L)},
\label{eq:pooling}
\end{equation}

where $V_G$ is the set of nodes in graph $G$. Alternative pooling strategies such as sum or max pooling are possible, but mean pooling is chosen for its stability across graphs of varying sizes.

\subsection{Deep SVDD Objective}
\label{sec:svdd}

Following the Deep \ac{SVDD} framework \citep{ruff2018deepsvdd}, we learn a hypersphere that encloses the normal data representations. The model maintains a center vector $\mathbf{c} \in \mathbb{R}^d$, which is initialized as the mean of the training embeddings:

\begin{equation}
\mathbf{c} = \frac{1}{|\mathcal{D}_{\text{train}}|} \sum_{\mathbf{z}_i \in \mathcal{D}_{\text{train}}} \mathbf{z}_i,
\label{eq:center_init}
\end{equation}

where $\mathcal{D}_{\text{train}}$ is the set of (node or graph) embeddings from the training data. The training objective minimizes the mean squared Euclidean distance from each embedding to the center:

\begin{equation}
\mathcal{L}_{\text{SVDD}} = \frac{1}{n} \sum_{i=1}^{n} \|\mathbf{z}_i - \mathbf{c}\|^2,
\label{eq:svdd_loss}
\end{equation}

with $n$ being the number of training samples. This loss encourages the encoder to map normal instances close to the center while implicitly pushing anomalies farther away.

\subsection{Anomaly Scoring}
\label{sec:scoring}

At inference time, the anomaly score of a sample is its squared distance to the center:

\begin{equation}
s(\mathbf{z}) = \|\mathbf{z} - \mathbf{c}\|^2.
\label{eq:anomaly_score}
\end{equation}

A higher score indicates a larger deviation from the normal prototype and thus a higher likelihood of being an anomaly. A decision boundary is typically defined as a quantile (e.g., the $95$th percentile) of the scores observed on a validation set.

\subsection{Node-Level versus Graph-Level}
\label{sec:node_vs_graph}

The \ac{OCGIN} architecture naturally supports both node-level and graph-level anomaly detection. In the node-level setting, each node $v \in V$ produces an embedding $\mathbf{z}_v$ and scores $s(\mathbf{z}_v)$ independently. In the graph-level setting, each graph $G$ produces a single embedding $\mathbf{z}_G$ via pooling, and scoring is performed per graph. Apart from the pooling step, the encoder architecture and loss function remain identical for both settings.


\section{Generating Logical Constraints with Decision Trees}
\label{sec:dt_constraints}

The core idea of this work is to extract interpretable logical rules from the data and use them as differentiable constraints to guide the OCGIN training. Decision trees provide a natural mechanism for this: they partition the feature space into axis-aligned regions, each described by a conjunction of simple threshold conditions. We train a decision tree on structural graph features and convert each root-to-leaf path into a logical rule.

\subsection{Feature Extraction}
\label{sec:feature_extraction}

To obtain features that capture structural properties of nodes and graphs, we compute a set of hand-crafted descriptors before training the decision tree.

\subsubsection{Node-Level Features}

For each node $v \in V$, we extract:

\begin{enumerate}
  \item \textbf{Raw node features:} The original attribute vector $\mathbf{x}_v \in \mathbb{R}^{d_{\text{in}}}$, where each dimension is treated as a separate feature. For citation networks, these typically represent bag-of-words encoding of the paper's content.
  \item \textbf{Node degree:} The number of incident edges:
  \begin{equation}
  \deg(v) = |\mathcal{N}(v)|.
  \label{eq:degree}
  \end{equation}
  \item \textbf{Local clustering coefficient:} The fraction of possible triangles through $v$ that actually exist:
  \begin{equation}
  \text{CC}(v) = \frac{2 \cdot |\{ (u,w) \in E : u,w \in \mathcal{N}(v) \}|}{\deg(v) \cdot (\deg(v) - 1)},
  \label{eq:clustering}
  \end{equation}
  where the numerator counts edges between neighbors of $v$. By convention, $\text{CC}(v) = 0$ if $\deg(v) < 2$.
\end{enumerate}

\subsubsection{Graph-Level Features}

For each graph $G$, we aggregate the node-level statistics:

\begin{enumerate}
  \item \textbf{Mean node features:} The per-dimension mean of all node feature vectors:
  \begin{equation}
  \bar{\mathbf{x}}_G^{(d)} = \frac{1}{|V_G|} \sum_{v \in V_G} \mathbf{x}_v^{(d)}, \quad d = 1, \dots, d_{\text{in}}.
  \label{eq:mean_features}
  \end{equation}
  \item \textbf{Mean node degree:} The average degree across all nodes in the graph:
  \begin{equation}
  \overline{\deg}(G) = \frac{1}{|V_G|} \sum_{v \in V_G} \deg(v).
  \label{eq:mean_degree}
  \end{equation}
\end{enumerate}

Each feature dimension is assigned a numeric index. The resulting feature vectors $\mathbf{X} \in \mathbb{R}^{m \times p}$, where $m$ is the number of samples (nodes or graphs) and $p$ is the total number of feature dimensions, serve as input to the decision tree.

\subsection{Decision Tree Induction}
\label{sec:tree_induction}

We train a binary decision tree that recursively partitions the feature space using the \ac{CART} algorithm \citep{breiman1984cart}. At each internal node, the algorithm searches for the feature $f$ and threshold $\theta$ that produce the purest split according to the Gini impurity.

\subsubsection{Gini Impurity}

For a set of samples $D$ with binary labels $y \in \{0, 1\}$, the Gini impurity measures class imbalance:

\begin{equation}
G(D) = 1 - \sum_{i \in \{0, 1\}} p_i^2,
\label{eq:gini}
\end{equation}

where $p_i = |\{y \in D : y = i\}| / |D|$ is the proportion of class $i$ in $D$. The impurity is zero when all samples belong to the same class and maximized at $0.5$ when classes are evenly split.

\subsubsection{Split Criterion}

Given a candidate split that partitions $D$ into left and right subsets $D_{\leq}$ and $D_{>}$ based on $x_f \leq \theta$, the weighted Gini impurity is:

\begin{equation}
G_{\text{split}}(f, \theta) = \frac{|D_{\leq}|}{|D|} \, G(D_{\leq}) + \frac{|D_{>}|}{|D|} \, G(D_{>}).
\label{eq:gini_split}
\end{equation}

The algorithm selects the split with minimum $G_{\text{split}}$:

\begin{equation}
(f^*, \theta^*) = \argmin_{f, \theta} G_{\text{split}}(f, \theta).
\label{eq:best_split}
\end{equation}

This process recurses on the two child partitions until one of the stopping criteria is met.

\subsubsection{Stopping Criteria}

A node is declared a leaf and no further split is attempted when:

\begin{enumerate}
  \item All samples in the node share the same label ($G(D) = 0$).
  \item The maximum tree depth (a hyperparameter) has been reached.
  \item No valid split exists that produces non-empty children (e.g., all feature values are identical).
\end{enumerate}

At a leaf node, the predicted class is the majority class of the samples in that leaf.

\subsection{Rule Extraction and Output Format}
\label{sec:rule_extraction}

Once the tree is trained, every root-to-leaf path defines a conjunctive rule. Each internal node along the path contributes an atomic condition comparing a feature to a threshold. A path of length $L$ produces the rule:

\begin{equation}
r_j: \bigwedge_{\ell = 1}^{L} \bigl( x_{f_\ell} \bowtie_\ell \theta_\ell \bigr) \implies \hat{y}_j,
\label{eq:rule_form}
\end{equation}

where $\bowtie_\ell \in \{\leq, >\}$ is the comparison operator and $\hat{y}_j$ is the majority class of the leaf. The conjunction of all conditions along the path must hold for a sample to be assigned to that leaf's class. The complete decision tree represents a disjunction of such conjunctive rules:

\begin{equation}
\hat{y}(\mathbf{x}) = \bigvee_{j: \mathbf{x} \text{ reaches leaf } j} r_j(\mathbf{x}).
\label{eq:dnf}
\end{equation}

Each rule is serialized into a JSON structure for downstream consumption:

\begin{verbatim}
{
  "conditions": [
    {"feature_index": 729, "op": "<=", "threshold": 0.0},
    {"feature_index": 1263, "op": "<=", "threshold": 0.0}
  ],
  "predicted_class": 1,
  "meta": {"depth": 2, "y_label_number": "20 v.s. 45"}
}
\end{verbatim}

The \texttt{feature\_index} refers to the column index in the feature matrix $\mathbf{X}$, \texttt{op} is the comparison operator, and \texttt{threshold} is the split value. The \texttt{predicted\_class} indicates whether samples satisfying this rule are predicted as normal ($0$) or anomalous ($1$). A separate mapping \texttt{additional\_attributes} records which semantic feature each index corresponds to, enabling traceability.

For the purpose of constraint generation, we retain only rules whose \texttt{predicted\_class} equals the normal label. These rules describe what constitutes normal structural behavior; violating them indicates potential anomaly.


\section{Differentiable Constraint Evaluation}
\label{sec:constraint_eval}

The logical rules extracted from the decision tree are discrete and non-differentiable. To use them within a gradient-based training loop, we convert each rule into a continuous value $L_c \in [0, 1]$ representing the degree to which a sample satisfies the anomalous rule set. We propose two alternative approaches: soft logic and distance-based scoring.

\subsection{Soft Logic-Based Handler}
\label{sec:soft_logic}

The soft logic approach replaces each hard comparison with a sigmoid approximation, making the entire rule evaluation differentiable.

\subsubsection{Atomic Condition Softening}

A condition $x_f \leq \theta$ is softened as:

\begin{equation}
\text{soft-leq}(x_f, \theta) = \sigma\bigl(\lambda (\theta - x_f)\bigr),
\label{eq:soft_leq}
\end{equation}

where $\sigma(\cdot)$ is the logistic sigmoid $\sigma(a) = 1 / (1 + e^{-a})$ and $\lambda > 0$ is a scaling factor controlling the sharpness of the approximation. As $\lambda \to \infty$, $\text{soft-leq}$ approaches the hard indicator $\mathbbm{1}[x_f \leq \theta]$. Conversely, $x_f > \theta$ is softened as:

\begin{equation}
\text{soft-gt}(x_f, \theta) = \sigma\bigl(\lambda (x_f - \theta)\bigr).
\label{eq:soft_gt}
\end{equation}

Note that $\text{soft-gt}(x_f, \theta) = 1 - \text{soft-leq}(x_f, \theta)$ only in the limit $\lambda \to \infty$; for finite $\lambda$, they form a soft partition.

\subsubsection{Logical Connectives}

Conjunction (AND) of conditions within a rule is implemented as the product of their soft values:

\begin{equation}
\text{soft-and}(v_1, \dots, v_k) = \prod_{i=1}^{k} v_i.
\label{eq:soft_and}
\end{equation}

This is the natural differentiable generalization of Boolean AND: the result is $1$ only if all $v_i$ are close to $1$, and $0$ if any $v_i$ is close to $0$.

Disjunction (OR) across multiple rules is implemented using De Morgan's law:

\begin{equation}
\text{soft-or}(v_1, \dots, v_m) = 1 - \prod_{i=1}^{m} (1 - v_i).
\label{eq:soft_or}
\end{equation}

This yields $0$ only if all $v_i$ are close to $0$, and $1$ if any $v_i$ is close to $1$.

\subsubsection{Constraint Value Computation}

For a set of anomaly rules $\mathcal{R}_{\text{anom}}$ (rules whose \texttt{predicted\_class} differs from the normal label), the constraint value for a sample $\mathbf{x}$ is:

\begin{equation}
L_c(\mathbf{x}) = \lambda_c \cdot \text{soft-or}\Bigl( \bigl\{ \text{soft-and}\bigl( \{ \text{cond}(c, \mathbf{x}) : c \in r \} \bigr) : r \in \mathcal{R}_{\text{anom}} \bigr\} \Bigr),
\label{eq:soft_constraint}
\end{equation}

where $\text{cond}(c, \mathbf{x})$ applies \eqref{eq:soft_leq} or \eqref{eq:soft_gt} depending on the operator in condition $c$, and $\lambda_c$ is a scaling factor. The resulting $L_c(\mathbf{x})$ lies in $[0, \lambda_c]$ and is high when $\mathbf{x}$ structurally resembles the training data (satisfies normal-class rules), and low when it deviates.

\subsection{Distance-Based Constraint Scoring}
\label{sec:distance_scoring}

The distance-based approach takes a different perspective: instead of measuring rule satisfaction directly, it computes the signed distance from a sample to the decision boundary of each rule. This provides a more granular signal, especially for samples that do not clearly satisfy or violate any rule.

\subsubsection{Distance to a Single Rule}

For a rule $r$ with conditions $\{(f_\ell, \bowtie_\ell, \theta_\ell)\}_{\ell=1}^L$, the distance to its decision boundary is defined as the minimum signed distance across all conditions:

\begin{equation}
d(\mathbf{x}, r) = \min_{\ell=1,\dots,L} \bigl( \mathbf{x}_{f_\ell} - \theta_\ell \bigr).
\label{eq:distance_rule}
\end{equation}

This quantity is positive when $\mathbf{x}$ satisfies the condition (is on the "normal side" of the boundary) and negative otherwise.

\subsubsection{Weighted Group Distance}

Rules are partitioned into two groups: normality rules $\mathcal{R}_{\text{norm}}$ (predicted class = normal) and anomaly rules $\mathcal{R}_{\text{anom}}$ (predicted class = anomaly). For each group, we compute the weighted mean distance:

\begin{equation}
\text{dist}(\mathbf{x}, \mathcal{R}) = \frac{\lambda_c}{|\mathcal{R}|} \sum_{r \in \mathcal{R}} d(\mathbf{x}, r),
\label{eq:group_distance}
\end{equation}

where $\lambda_c$ is a scaling factor controlling the influence of the constraint.

\subsubsection{Hard Rule Satisfaction}

In addition to the continuous distance, we evaluate hard satisfaction to determine which region $\mathbf{x}$ falls into:

\begin{equation}
\text{satisfies}(\mathbf{x}, r) = \begin{cases}
\text{True} & \text{if } \forall \ell: \mathbf{x}_{f_\ell} \bowtie_\ell \theta_\ell, \\
\text{False} & \text{otherwise}.
\end{cases}
\label{eq:hard_satisfaction}
\end{equation}

\subsubsection{Three-Zone Scoring}

The final unnormalized score is computed based on three mutually exclusive cases:

\begin{enumerate}
  \item \textbf{Normal region:} $\mathbf{x}$ satisfies at least one normality rule. In this case, the constraint score combines the distance to the satisfied normality boundary with the minimum distance to any anomaly rule:
  \begin{equation}
  \text{score}(\mathbf{x}) = -\, d(\mathbf{x}, r_{\text{norm}}) \cdot \min_{r \in \mathcal{R}_{\text{anom}}} d(\mathbf{x}, r) + \bigl( \text{dist}(\mathbf{x}, \mathcal{R}_{\text{norm}}) - \text{dist}(\mathbf{x}, \mathcal{R}_{\text{anom}}) \bigr).
  \label{eq:normal_region}
  \end{equation}
  
  \item \textbf{Anomaly region:} $\mathbf{x}$ satisfies at least one anomaly rule:
  \begin{equation}
  \text{score}(\mathbf{x}) = d(\mathbf{x}, r_{\text{anom}}) \cdot \min_{r \in \mathcal{R}_{\text{norm}}} d(\mathbf{x}, r) + \bigl( \text{dist}(\mathbf{x}, \mathcal{R}_{\text{norm}}) - \text{dist}(\mathbf{x}, \mathcal{R}_{\text{anom}}) \bigr).
  \label{eq:anomaly_region}
  \end{equation}
  
  \item \textbf{Grey zone:} $\mathbf{x}$ satisfies neither group. Only the group distance difference contributes:
  \begin{equation}
  \text{score}(\mathbf{x}) = \text{dist}(\mathbf{x}, \mathcal{R}_{\text{norm}}) - \text{dist}(\mathbf{x}, \mathcal{R}_{\text{anom}}).
  \label{eq:grey_zone}
  \end{equation}
\end{enumerate}

The final constraint value is obtained by applying the sigmoid function:

\begin{equation}
L_c(\mathbf{x}) = \sigma\bigl(\text{score}(\mathbf{x})\bigr) = \frac{1}{1 + e^{-\text{score}(\mathbf{x})}}.
\label{eq:dist_final}
\end{equation}

This maps the score to the $(0, 1)$ interval, where values near $1$ indicate anomaly-like structural patterns (high constraint activation) and values near $0$ indicate normal-like patterns.


\section{Integrating Constraints into OCGIN}
\label{sec:integration}

The constraint value $L_c \in [0, 1]$ produced by either the soft logic or distance-based handler is incorporated into the OCGIN framework at two possible points: the loss function during training and the anomaly score during inference.

\subsection{Constraint-Augmented Loss Functions}
\label{sec:constrained_loss}

We define three variants for integrating $L_c$ into the Deep SVDD loss. In all cases, $\mathcal{L}_{\text{SVDD}}$ is defined as in \eqref{eq:svdd_loss} and $L_c^{(i)}$ denotes the constraint value for the $i$-th sample.

\subsubsection{Additive Constraint Loss}

The constraint term is added directly to the SVDD loss:

\begin{equation}
\mathcal{L}_{\text{add}} = \mathcal{L}_{\text{SVDD}} + \frac{1}{n} \sum_{i=1}^{n} L_c^{(i)}.
\label{eq:loss_add}
\end{equation}

This formulation treats high constraint values as an additional penalty: samples that structurally resemble the training data increase the overall loss, pushing the model to learn tighter representations for them.

\subsubsection{Weighted Constraint Loss}

The SVDD loss per sample is scaled by $(1 + L_c)$:

\begin{equation}
\mathcal{L}_{\text{weight}} = \frac{1}{n} \sum_{i=1}^{n} \| \mathbf{z}_i - \mathbf{c} \|^2 \cdot \bigl(1 + L_c^{(i)}\bigr).
\label{eq:loss_weight}
\end{equation}

Samples with high constraint values contribute more to the total loss, effectively emphasizing structurally normal patterns during training.

\subsubsection{Ignore-Suspicious Loss}

The SVDD loss per sample is scaled by $(1 - L_c)$:

\begin{equation}
\mathcal{L}_{\text{ignore}} = \frac{1}{n} \sum_{i=1}^{n} \| \mathbf{z}_i - \mathbf{c} \|^2 \cdot \bigl(1 - L_c^{(i)}\bigr).
\label{eq:loss_ignore}
\end{equation}

This has the opposite effect: samples with high constraint values (structurally normal according to the decision tree) are down-weighted. This formulation is intended for scenarios where the decision tree's normal-class rules may be unreliable and should be treated with caution.

\subsection{Constraint-Aware Anomaly Scoring}
\label{sec:constrained_scoring}

Beyond influencing training, constraints can also be incorporated into the anomaly score at inference time. This is useful when we want the decision tree's structural knowledge to persist beyond the training phase. We extend \eqref{eq:anomaly_score} to include $L_c$:

\begin{equation}
s_{\text{final}}(\mathbf{z}_i, L_c^{(i)}) = \begin{cases}
\|\mathbf{z}_i - \mathbf{c}\|^2 + L_c^{(i)} & \text{(additive)}, \\[4pt]
\|\mathbf{z}_i - \mathbf{c}\|^2 \cdot \bigl(1 + L_c^{(i)}\bigr) & \text{(weighting)}, \\[4pt]
\|\mathbf{z}_i - \mathbf{c}\|^2 \cdot \bigl(1 - L_c^{(i)}\bigr) & \text{(ignore-suspicious)}.
\end{cases}
\label{eq:forecast_score}
\end{equation}

The three modes correspond directly to the loss variants. When this forecasting extension is used, constraints influence both the learned representation (through the loss) and the final decision (through the score). When it is not used, only the anomaly score $s(\mathbf{z}) = \|\mathbf{z} - \mathbf{c}\|^2$ is employed at test time, and constraints solely affect training.

\subsection{Training Procedure}
\label{sec:training_procedure}

The overall training procedure proceeds as follows:

\begin{enumerate}
  \item Split the training data into a main training set and a separate tree-induction set.
  \item Train the decision tree on the tree-induction set using the features described in Section~\ref{sec:feature_extraction}.
  \item Extract rules and serialize them to JSON.
  \item Initialize the constraint handler (soft logic or distance-based) with the rule file.
  \item Compute $L_c$ for all training samples using the chosen handler.
  \item Initialize the OCGIN center $\mathbf{c}$ on the main training set \eqref{eq:center_init}.
  \item Train the OCGIN encoder by minimizing one of the constrained losses \eqref{eq:loss_add}--\eqref{eq:loss_ignore}, where $L_c$ is precomputed and fixed throughout training.
  \item At test time, compute the final anomaly score using \eqref{eq:anomaly_score} or \eqref{eq:forecast_score} as appropriate.
\end{enumerate}

The decision tree is trained once and kept fixed during OCGIN training. This design ensures that the constraints provide a stable, interpretable inductive bias throughout the optimization process.
